Beide überwinden dieselbe Höhe $h$. Die Hubarbeit $W = m\,g\,h$ ist proportional zur Masse:
\begin{align}
 \frac{W_2-W_1}{W_1} &= \frac{m_2\,g\,h-m_1\,g\,h}{m_1\,g\,h} = \relF \\
   &= \frac{\mpb-\mpa}{\mpa} = \rel \approx \result{\relP{0}{2}}
\end{align}
Da die beiden gemeinsam gehen, brauchen sie dieselbe Zeit $t$. Die Leistung $P = W/t$ ist deshalb ebenfalls um $\result{\relP{0}{2}}$ grösser.
